% Above: how an event gets an honest timestamp. Below: the protocol that
% design makes possible, from the 1994 paper. Nothing below is built until
% chapter 12. Notes sit under each half rather than beside it, so that they
% cannot land on a lifeline.
\begin{tikzpicture}[font=\sffamily\scriptsize, x=1cm, y=1cm]

% ================================================================ upper half
\node[chlabel, anchor=west] at (0,0.95) {stamp at the event, not at the read};
\foreach \x/\name/\id in {0/{sensor pin}/pin, 3.4/{TIM2 capture}/cap,
                          6.8/{capture handler}/isr, 10.2/{node state}/st}{
  \node[actor, text width=20mm] (\id) at (\x,0) {\name};
  \draw[lifeline] (\id.south) -- ++(0,-4.1);
}
\draw[msg] (0,-1.0) -- node[above]{edge} (3.4,-1.0);
\node[activation, minimum height=3mm] at (3.4,-1.2) {};
\draw[msg] (3.4,-1.9) -- node[above]{interrupt, possibly late} (6.8,-1.9);
\node[activation, minimum height=4mm] at (6.8,-2.15) {};
\draw[reply] (6.8,-2.65) -- node[below]{read CCR1: the count at the edge} (3.4,-2.65);
\draw[msg] (6.8,-3.7) -- node[above]{sample plus its true instant} (10.2,-3.7);

\node[umlnote, text width=104mm, anchor=north west] at (-0.6,-4.5)
  {Interrupt latency moves when the node \emph{learns} about the event, not when
   the event is \emph{recorded}. The same sample stamped when a two-wire read
   completes would instead carry the whole bus transaction, the queueing, and
   whatever else the node was doing at the time. Chapter 6 depends on this
   entirely.};

% ================================================================ lower half
\node[chlabel, anchor=west] at (0,-6.5)
  {the protocol this makes possible, designed here and built in chapter 12};
\foreach \x/\name/\id in {0/{node A, the marker}/a, 4.6/{the bus}/bus, 9.2/{node B}/b}{
  \node[actor, text width=24mm] (\id) at (\x,-7.35) {\name};
  \draw[lifeline] (\id.south) -- ++(0,-4.6);
}

\draw[msg] (0,-8.0) -- node[above]{frame A, "mark", no payload} (4.6,-8.0);
\node[activation, minimum height=3mm] at (0,-8.2) {};
\node[tlabel, anchor=west] at (0.15,-8.75) {A stamps it too};

\draw[msg] (4.6,-8.7) -- node[above]{every node stamps reception} (9.2,-8.7);
\node[activation, minimum height=3mm] at (9.2,-8.9) {};

\draw[msg] (0,-9.7) -- node[above]{frame B, follow-up: A's own stamp of frame A} (9.2,-9.7);

\draw[msg] (9.2,-10.5) -- ++(1.0,0) -- ++(0,-0.5)
  node[right, xshift=1mm]{offset = A's stamp minus mine, then slew} -- ++(-1.0,0);

\node[umlnote, text width=104mm, anchor=north west] at (-0.6,-11.9)
  {The time-critical path is reception to timestamp, which is local to each node
   and can be characterised once. On this part the bus controller removes even
   that, by capturing at the start-of-frame bit before any interrupt runs. Under
   twenty messages a second and one identifier, for about twenty microseconds in
   the 1994 result and below one microsecond in the 2021 one, the difference
   between them being rate correction rather than offset alone.};
\end{tikzpicture}
